\documentclass[11pt]{article}
\usepackage{amsmath,amsthm,amssymb}
\usepackage[margin=1in]{geometry}

\newtheorem{theorem}{Theorem}
\newtheorem{lemma}[theorem]{Lemma}
\newtheorem{proposition}[theorem]{Proposition}
\newtheorem{corollary}[theorem]{Corollary}
\theoremstyle{definition}
\newtheorem{definition}[theorem]{Definition}
\theoremstyle{remark}
\newtheorem{remark}[theorem]{Remark}

\DeclareMathOperator{\id}{id}

\title{The Cactus Midpoint Theorem}
\author{Clio}
\date{12 April 2026}

\begin{document}
\maketitle

\begin{abstract}
We prove that in the 0-Hecke algebra $H_0(S_k)$, the Tikhonov interpolating
operators $R_i(p) = pT_i + (1-p)$ evaluated at $p = 1/2$ yield a scalar
multiple of the cactus (Demazure) sorting operator:
$\prod R_{i_j}(1/2) = (1/2)^m \pi_{w_0^{[a,b]}}$. We show this identity
is \emph{special to $q = 0$} by proving that the Demazure elements
$\pi_i = T_i + 1$ in $H_q(S_k)$ satisfy a braid obstruction
$\pi_i\pi_{i+1}\pi_i - \pi_{i+1}\pi_i\pi_{i+1} = q(T_i - T_{i+1})$,
which vanishes precisely at $q = 0$. We also identify the Tikhonov spectral
parameter $\kappa = 2 - q$ at the cactus midpoint.
\end{abstract}

\section{Setup}

Let $H_q(S_k)$ be the Iwahori--Hecke algebra of type $A_{k-1}$ over
$\mathbb{Q}(q)$, with generators $T_1, \ldots, T_{k-1}$ satisfying:
\begin{align}
T_i^2 &= (q-1)T_i + q, \label{eq:hecke-quadratic} \\
T_i T_{i+1} T_i &= T_{i+1} T_i T_{i+1}, \label{eq:braid} \\
T_i T_j &= T_j T_i \quad \text{for } |i-j| \geq 2. \label{eq:commute}
\end{align}

\begin{definition}
The \emph{Demazure element} (or $\pi$-operator) is $\pi_i := T_i + 1$.
\end{definition}

\begin{definition}
The \emph{Tikhonov interpolating operator} is
$R_i(p) := p T_i + (1-p) \in H_q(S_k)[p]$.
\end{definition}

For an interval $[a,b] \subseteq [1,k]$, let $w_0^{[a,b]}$ denote the
longest element of the parabolic subgroup $S_{[a,b]}$, with
$m = \ell(w_0^{[a,b]}) = \binom{b-a+1}{2}$.

\section{The Demazure Algebra}

\begin{lemma}[Quasi-idempotency]\label{lem:quasi-idem}
In $H_q(S_k)$, the Demazure elements satisfy
\[
\pi_i^2 = (q+1)\,\pi_i.
\]
In particular, $\pi_i$ is idempotent if and only if $q = 0$.
\end{lemma}

\begin{proof}
Expanding:
\begin{align*}
\pi_i^2 = (T_i + 1)^2 &= T_i^2 + 2T_i + 1 \\
&= \bigl((q-1)T_i + q\bigr) + 2T_i + 1 \\
&= (q+1)T_i + (q+1) \\
&= (q+1)(T_i + 1) = (q+1)\,\pi_i.
\end{align*}
At $q = 0$ this gives $\pi_i^2 = \pi_i$.
\end{proof}

\begin{theorem}[Braid Obstruction]\label{thm:braid-obstruction}
For adjacent generators $\pi_i, \pi_{i+1}$ in $H_q(S_k)$:
\begin{equation}\label{eq:braid-obs}
\pi_i \pi_{i+1} \pi_i - \pi_{i+1} \pi_i \pi_{i+1} = q(T_i - T_{i+1}).
\end{equation}
The braid relation for $\pi_i$ holds if and only if $q = 0$.
\end{theorem}

\begin{proof}
We expand both sides using the Hecke relations.

\medskip\noindent\textbf{Left side.}
First compute $(T_{i+1}+1)(T_i+1) = T_{i+1}T_i + T_{i+1} + T_i + 1$.
Then:
\begin{align*}
\pi_i\pi_{i+1}\pi_i
&= (T_i + 1)\bigl(T_{i+1}T_i + T_{i+1} + T_i + 1\bigr) \\
&= T_i T_{i+1} T_i + T_i T_{i+1} + T_i^2 + T_i
   + T_{i+1}T_i + T_{i+1} + T_i + 1 \\
&= T_i T_{i+1} T_i + T_i T_{i+1} + T_{i+1}T_i + T_{i+1} + (q+1)T_i + (q+1),
\end{align*}
where we used $T_i^2 = (q-1)T_i + q$ to get the last line:
$T_i^2 + T_i = (q-1)T_i + q + T_i = qT_i + q$, and the remaining $T_i$
gives $(q+1)T_i + (q+1)$ total.

\medskip\noindent\textbf{Right side.}
By symmetry under $i \leftrightarrow i+1$:
\[
\pi_{i+1}\pi_i\pi_{i+1}
= T_{i+1}T_i T_{i+1} + T_{i+1}T_i + T_i T_{i+1} + T_i + (q+1)T_{i+1} + (q+1).
\]

\medskip\noindent\textbf{Difference.}
Using the braid relation $T_iT_{i+1}T_i = T_{i+1}T_iT_{i+1}$, the
leading terms cancel:
\begin{align*}
\pi_i\pi_{i+1}\pi_i - \pi_{i+1}\pi_i\pi_{i+1}
&= \bigl[T_{i+1} + (q+1)T_i\bigr] - \bigl[T_i + (q+1)T_{i+1}\bigr] \\
&= qT_i - qT_{i+1} \\
&= q(T_i - T_{i+1}).
\end{align*}
The terms $T_iT_{i+1} + T_{i+1}T_i$ appear identically on both sides.
\end{proof}

\begin{corollary}[Word-Independence at $q = 0$]\label{cor:word-indep}
At $q = 0$, for any $w \in S_k$, the product $\pi_w := \prod_{j=1}^m \pi_{i_j}$
over a reduced word $s_{i_1}\cdots s_{i_m}$ of $w$ is independent of the
choice of reduced word.
\end{corollary}

\begin{proof}
The reduced words for $w$ are connected by braid moves
($s_i s_{i+1} s_i \leftrightarrow s_{i+1} s_i s_{i+1}$) and commutation
moves ($s_i s_j \leftrightarrow s_j s_i$ for $|i-j| \geq 2$) by Matsumoto's
theorem. By Theorem~\ref{thm:braid-obstruction}, the braid relation
$\pi_i\pi_{i+1}\pi_i = \pi_{i+1}\pi_i\pi_{i+1}$ holds at $q = 0$.
The commutation relation $\pi_i\pi_j = \pi_j\pi_i$ for $|i-j| \geq 2$
follows from \eqref{eq:commute}. Hence the product is word-independent.
\end{proof}

\begin{corollary}[Word-Dependence at Generic $q$]\label{cor:word-dep}
At generic $q \neq 0$, the product $\prod \pi_{i_j}$ over a reduced word
for $w_0$ depends on the choice of reduced word. In the regular representation
of $H_q(S_3)$, the product $(T_1+1)(T_2+1)(T_1+1)$ has rank 3,
not rank 1. It is not a cactus commutor.
\end{corollary}

\begin{proof}
The obstruction $q(T_1 - T_2) \neq 0$ for $q \neq 0$ shows word-dependence.
For the rank claim, the product acts on the 6-dimensional regular
representation. At $q = 0$, idempotency forces the image of each $\pi_i$ to
be a proper subspace, and the composition collapses to rank 1. At $q = 1$
(the symmetric group), direct computation shows rank 3
(verified computationally in $H_1(S_3)$). Since the rank is an upper
semicontinuous function of $q$, the rank is $\geq 3$ at generic $q$.
\end{proof}

\section{The Cactus Midpoint Theorem}

\begin{lemma}[Factorization of $R_i(p)$]\label{lem:factorization}
The Tikhonov operator factors as:
\begin{equation}\label{eq:Rip-factor}
R_i(p) = p\,\pi_i + (1-2p)\cdot\id.
\end{equation}
\end{lemma}

\begin{proof}
$R_i(p) = pT_i + (1-p) = p(T_i + 1) - p + (1-p) = p\,\pi_i + (1-2p)$.
\end{proof}

\begin{theorem}[Cactus Midpoint]\label{thm:cactus-midpoint}
At $q = 0$, for any interval $[a,b] \subseteq [1,k]$ and any reduced word
$s_{i_1}\cdots s_{i_m}$ of $w_0^{[a,b]}$:
\begin{equation}\label{eq:midpoint}
\prod_{j=1}^{m} R_{i_j}\!\left(\tfrac{1}{2}\right) = \left(\tfrac{1}{2}\right)^m \pi_{w_0^{[a,b]}}.
\end{equation}
\end{theorem}

\begin{proof}
By Lemma~\ref{lem:factorization}, at $p = 1/2$:
\[
R_i\!\left(\tfrac{1}{2}\right) = \tfrac{1}{2}\,\pi_i + 0 = \tfrac{1}{2}\,\pi_i.
\]
The identity term $(1-2p)$ vanishes at $p = 1/2$. Therefore:
\[
\prod_{j=1}^{m} R_{i_j}\!\left(\tfrac{1}{2}\right)
= \prod_{j=1}^{m} \tfrac{1}{2}\,\pi_{i_j}
= \left(\tfrac{1}{2}\right)^m \prod_{j=1}^{m} \pi_{i_j}
= \left(\tfrac{1}{2}\right)^m \pi_{w_0^{[a,b]}},
\]
where the last equality uses word-independence
(Corollary~\ref{cor:word-indep}) and the definition
$\pi_{w_0^{[a,b]}} = \prod \pi_{i_j}$.
\end{proof}

\begin{theorem}[M\"obius Factorization]\label{thm:mobius}
For the product over a reduced word of $w_0^{[a,b]}$ at $q = 0$:
\begin{equation}\label{eq:mobius}
\prod_{j=1}^m R_{i_j}(p) = \sum_{S \subseteq \{1,\ldots,m\}} p^{|S|}(1-2p)^{m-|S|} \cdot \vec{\pi}_S,
\end{equation}
where $\vec{\pi}_S = \prod_{j \in S}^{\to} \pi_{i_j}$ is the ordered
product of $\pi_{i_j}$ for $j \in S$ (in the order inherited from the
reduced word), and identity operators fill positions $j \notin S$.
At $p = 1/2$, only the term $S = \{1,\ldots,m\}$ survives, recovering
Theorem~\ref{thm:cactus-midpoint}.
\end{theorem}

\begin{proof}
By Lemma~\ref{lem:factorization}, $R_{i_j}(p) = p\,\pi_{i_j} + (1-2p)\cdot\id$.
Expanding the product:
\[
\prod_{j=1}^m \bigl(p\,\pi_{i_j} + (1-2p)\cdot\id\bigr)
= \sum_{S \subseteq \{1,\ldots,m\}} \prod_{j=1}^m
\begin{cases} p\,\pi_{i_j} & j \in S \\ (1-2p)\cdot\id & j \notin S \end{cases}
= \sum_S p^{|S|}(1-2p)^{m-|S|}\,\vec{\pi}_S.
\]
At $p = 1/2$, $(1-2p) = 0$, so only $S = \{1,\ldots,m\}$ contributes:
$p^m \vec{\pi}_{\{1,\ldots,m\}} = (1/2)^m\,\pi_{w_0^{[a,b]}}$.
\end{proof}

\section{The Eigenvalue Picture}

\begin{proposition}[Eigenvalues of $R_i(p)$]\label{prop:eigenvalues}
The generator $T_i$ has eigenvalues $q$ (symmetric eigenspace) and $-1$
(antisymmetric eigenspace). The operator $R_i(p) = pT_i + (1-p)$ has
eigenvalues:
\begin{align}
\lambda_+ &= 1 - p(1-q) = \frac{1+q}{2}\Big|_{p=1/2}
  & &\text{on the $q$-eigenspace of $T_i$}, \label{eq:eig-plus}\\
\lambda_- &= 1 - 2p = 0\Big|_{p=1/2}
  & &\text{on the $(-1)$-eigenspace of $T_i$}. \label{eq:eig-minus}
\end{align}
At $p = 1/2$, $R_i(1/2)$ annihilates the antisymmetric eigenspace for
\emph{all} values of $q$, and scales the symmetric eigenspace by $(1+q)/2$.
\end{proposition}

\begin{proof}
If $T_i v = qv$, then $R_i(p)v = (pq + 1 - p)v = (1 - p(1-q))v$.
If $T_i v = -v$, then $R_i(p)v = (-p + 1 - p)v = (1-2p)v$.
Setting $p = 1/2$ gives the stated values.
\end{proof}

\begin{remark}
The cactus midpoint $p = 1/2$ is the unique value of $p$ that kills the
antisymmetric eigenspace of \emph{every} $R_i$, independent of $q$. However,
the product $\prod R_{i_j}(1/2)$ fails to be word-independent at generic
$q$ due to the braid obstruction (Theorem~\ref{thm:braid-obstruction}).
Each individual $R_i(1/2)$ is a scaled projector, but their composition
does not give a well-defined element of the Hecke algebra when $q \neq 0$.
\end{remark}

\section{The Spectral Parameter}

\begin{proposition}[Tikhonov's Spectral Parameter]\label{prop:kappa}
Define $\kappa = \kappa(p,q) := \frac{1-qp}{1-p}$. Then:
\begin{enumerate}
\item[(a)] The random sorting permuton $\mu_{p,q}$ depends on $(p,q)$
  only through $\kappa$ (Tikhonov, Theorem 1.4 of \cite{tikhonov}).
\item[(b)] The three distinguished loci are:
\begin{center}
\begin{tabular}{lccc}
\hline
Level & $p$ & $\kappa$ & Operator \\
\hline
Identity & $0$ & $1$ & $\id$ \\
Cactus & $1/2$ & $2 - q$ & $(1/2)^m\,\pi_{w_0}$ \\
Hecke & $1$ & $\frac{1}{1-q}$ & $T_{w_0}$ \\
\hline
\end{tabular}
\end{center}
\item[(c)] At $q = 0$, the cactus level is $\kappa = 2$, independent of
  all other parameters. This is a universal constant: the midpoint between
  $\kappa = 1$ (identity) and $\kappa = \infty$ (fully sorted), measured
  on the scale $\kappa/(1+\kappa)$ which maps $[1,\infty] \to [1/2, 1]$.
\end{enumerate}
\end{proposition}

\begin{proof}
Part (a) is Tikhonov's theorem. For (b), direct computation:
at $p = 0$: $\kappa = 1/1 = 1$;
at $p = 1/2$: $\kappa = (1 - q/2)/(1/2) = 2 - q$;
at $p = 1$: $\kappa = (1-q)/0$, interpreted as the limit $\lim_{p \to 1} \kappa = +\infty$ when $q \neq 1$, or via $1/(1-q)$ in the algebraic sense.

For (c), at $q = 0$, $\kappa(1/2, 0) = 2$.
\end{proof}

\section{Why $q = 0$ is Special}

We collect the algebraic reasons why the cactus midpoint theorem is
special to $q = 0$.

\begin{theorem}\label{thm:why-q0}
The following properties hold at $q = 0$ and fail at generic $q$:
\begin{enumerate}
\item \textbf{Idempotency}: $\pi_i^2 = \pi_i$ (Lemma~\ref{lem:quasi-idem}).
\item \textbf{Braid relation}: $\pi_i\pi_{i+1}\pi_i = \pi_{i+1}\pi_i\pi_{i+1}$
  (Theorem~\ref{thm:braid-obstruction}).
\item \textbf{Word-independence}: $\pi_w$ is well-defined
  (Corollary~\ref{cor:word-indep}).
\item \textbf{Rank collapse}: The image of $\pi_{w_0}$ acting on the
  regular representation has dimension $1$ (all of $H_0(S_k)$ maps onto
  the span of $\sum_{v \leq w_0} T_v$).
\end{enumerate}
The chain of implications is: $(1) + (2) \Rightarrow (3) \Rightarrow (4)$.
\end{theorem}

\begin{proof}
Properties (1) and (2) are proved in Lemma~\ref{lem:quasi-idem} and
Theorem~\ref{thm:braid-obstruction}. Property (3) follows from (1), (2),
and Matsumoto's theorem as in Corollary~\ref{cor:word-indep}.

For (4): At $q = 0$, $\pi_i$ is an idempotent, hence its image is a proper
subspace. In the regular representation on basis $\{T_w : w \in S_k\}$, the
operator $\pi_i = T_i + 1$ maps $T_w$ to $T_{s_iw} + T_w$ if $\ell(s_iw) > \ell(w)$,
and to $0$ if $\ell(s_iw) < \ell(w)$ (since $T_i T_w = -T_w$ in $H_0$ when
$\ell(s_iw) < \ell(w)$). The product $\pi_{w_0}$ maps every basis element
to $\sum_{v \in S_k} T_v$, giving rank 1.

To see this, note that $\pi_{w_0}$ is the unique maximal element in the
$0$-Hecke monoid: applying all sorting operators in sequence sorts any
permutation to the identity, and the image of any $T_w$ under $\pi_{w_0}$
is $\sum_{v} T_v$ (the ``fully sorted'' element).

At generic $q \neq 0$, property (1) fails ($\pi_i^2 = (q+1)\pi_i \neq \pi_i$),
(2) fails (obstruction $q(T_i - T_{i+1}) \neq 0$), and in $H_q(S_3)$,
the product $(T_1+1)(T_2+1)(T_1+1)$ has rank 3 in the regular representation
(verified computationally at $q = 1$).
\end{proof}

\section{Computational Verification}

All results have been verified computationally:

\begin{enumerate}
\item \textbf{$H_q(S_3)$ regular representation} (6-dimensional):
  Hecke relations, quasi-idempotency, braid obstruction $q(T_i - T_{i+1})$,
  word-dependence at generic $q$, word-independence at $q = 0$.

\item \textbf{$H_0(S_3)$ and $H_0(S_4)$}: Cactus midpoint identity
  $\prod R_{i_j}(1/2) = (1/2)^m \pi_{w_0}$ verified for all parabolic
  sub-products.

\item \textbf{Rank analysis}:
  $\operatorname{rank}(\pi_1\pi_2\pi_1) = 1$ at $q = 0$,
  $= 3$ at $q = 1$, in $H_q(S_3)$.

\item \textbf{Spectral parameter}: $\kappa(1/2, q) = 2 - q$ confirmed
  symbolically.
\end{enumerate}

Scripts: \texttt{braid\_obstruction\_verify.py},
\texttt{tikhonov\_cactus\_test.py},
\texttt{cactus\_midpoint\_generic\_q.py}.

\begin{thebibliography}{10}
\bibitem{tikhonov} K.~Tikhonov, \emph{Random sorting networks and the
  Hecke algebra}, arXiv:2604.06532, 2026.
\end{thebibliography}

\end{document}
